\section{Lecture 14 (February 14, 2025, est.): Relative Tannaka duality}
\label{lec14}
\noindent\textit{Scribe: Guillermo Sanmarco}

\subsection{Changing the target of a fiber functor}

The next goal is to study symmetric tensor categories by realizing them as
categories of representations of affine group schemes in a suitable symmetric
tensor category. Ordinary Tannaka duality uses $\Vect_k$ as the target; the
relative version allows a different target $\mathcal D$.

Fix a perfect field $k$. In this lecture, a symmetric tensor category is
essentially small, $k$-linear, abelian and rigid, with finite-dimensional
morphism spaces, objects of finite length, and
$\End(\1)=k$. A tensor functor is exact, $k$-linear and strong symmetric
monoidal. Such a functor is faithful. Write $\operatorname{Symm}_k$ for the
category of these categories and isomorphism classes of tensor functors.
Its full subcategory of categories of moderate growth consists of those
$\mathcal C$ for which, for each $X\in\mathcal C$, there is a real number
$a_X\geq1$ satisfying
\[
 \operatorname{length}(X^{\otimes n})\leq a_X^n\qquad(n\geq1).
\]

Fix $\mathcal D\in\operatorname{Symm}_k$. An object of
$\operatorname{Symm}(\mathcal D)$ is a pair $(\mathcal C,F)$ with
$F:\mathcal C\to\mathcal D$ a tensor functor. A morphism from
$(\mathcal C_1,F_1)$ to $(\mathcal C_2,F_2)$ is represented by a pair
$(T,\alpha)$, where $T:\mathcal C_1\to\mathcal C_2$ is a tensor functor
and $\alpha:F_1\xrightarrow{\sim}F_2T$ is a monoidal natural isomorphism.
We identify two such pairs when there is a compatible monoidal natural
isomorphism between their functors $T$. Because $F_2$ is faithful, a
compatible isomorphism between two fixed pairs is unique if it exists.
This does not say that the morphisms of $\operatorname{Symm}(\mathcal D)$
themselves are unique.

\subsection{The internal group of tensor automorphisms}

An affine group scheme in $\mathcal D$ is described by a commutative Hopf
algebra in $\Ind(\mathcal D)$. The ind-completion is necessary: a coordinate
algebra need not be an object of finite length. If $F:\mathcal C\to\mathcal D$
is a tensor functor, the coefficient construction gives a commutative Hopf
algebra
\[
 A_F=\int^{X\in\mathcal C}F(X)^*\otimes F(X)
       \quad\text{in }\Ind(\mathcal D).
\]
Here the integral denotes a coend. Concretely, its structure maps
$i_X:F(X)^*\otimes F(X)\to A_F$ satisfy, for every $f:X\to Y$,
\[
 i_X\bigl(F(f)^*\otimes\id_{F(X)}\bigr)
 =i_Y\bigl(\id_{F(Y)^*}\otimes F(f)\bigr)
 \quad\text{on }F(Y)^*\otimes F(X).
\]
It is universal for families with these relations. Tensor products in
$\mathcal C$ define multiplication, the unit object defines the unit,
evaluation and coevaluation define the coalgebra structure, and duality
defines the antipode. Symmetry makes the multiplication commutative.
This is the internal analogue of the matrix-coefficient Hopf algebra.

\begin{definition}
The affine group scheme represented by $A_F$ is denoted
$\underline{\Aut}^{\otimes}(F)$. The \emph{fundamental group} of
$\mathcal D$ is
\[
 \pi(\mathcal D)=\underline{\Aut}^{\otimes}(\Id_{\mathcal D}).
\]
\end{definition}

The universal coefficient coactions make every object of $\mathcal D$, and
then every ind-object, a representation of $\pi(\mathcal D)$. Naturality
means that every morphism in $\mathcal D$ respects these canonical actions.
Applying the canonical action to the objects $F(X)$ gives a homomorphism
\[
 \psi_F:\pi(\mathcal D)\longrightarrow
             \underline{\Aut}^{\otimes}(F).
\]

\begin{example}
We have $\pi(\Vect_k)=1$. Indeed, naturality with respect to maps
$k\to V$, together with the tensor-unit condition, forces a tensor
automorphism of the identity to be the identity on every vector space.
If $\operatorname{char}k\ne2$, then
$\pi(\SVec_k)=\mu_2\cong\mathbb Z/2$, regarded as an internal group
concentrated in even degree. Its nontrivial element acts as $+1$ on the
even part and $-1$ on the odd part. This is the parity automorphism,
not the functor that exchanges even and odd degrees.
\end{example}

\subsection{Compatible groups and representations}

Not every homomorphism from $\pi(\mathcal D)$ to an internal group is
compatible with the structure already present in $\mathcal D$.

\begin{definition}
A \emph{$\mathcal D$-group} is a pair $(G,\phi)$, where $G$ is an affine
group scheme in $\mathcal D$ and $\phi:\pi(\mathcal D)\to G$ is a
homomorphism such that the canonical action of $\pi(\mathcal D)$ on
$\mathcal O(G)$ equals the action induced by conjugation in $G$ along
$\phi$.
A morphism $(G,\phi)\to(H,\psi)$ is a homomorphism $u:G\to H$ with
$u\phi=\psi$.
\end{definition}

For clarity, put $A=\mathcal O(G)$ and $B=\mathcal O(\pi(\mathcal D))$.
Using left coactions, the compatibility condition says that the canonical
map $A\to B\otimes A$ is the composite
\[
 A\xrightarrow{\Delta^{(2)}} A\otimes A\otimes A
 \xrightarrow{\phi^*\otimes\id\otimes S_B\phi^*}
 B\otimes A\otimes B
 \xrightarrow{(m_B\otimes\id)(\id\otimes c_{A,B})}
 B\otimes A.
\]
The last arrow uses the symmetry to bring the two $B$ factors together.
The pair $(\underline{\Aut}^{\otimes}(F),\psi_F)$ is a
$\mathcal D$-group.

\begin{definition}
For a $\mathcal D$-group $(G,\phi)$, let
$\Rep_{\mathcal D}(G,\phi)$ be the full subcategory of representations
of $G$ on objects of $\mathcal D$ whose restriction along $\phi$ is the
canonical $\pi(\mathcal D)$-action. It has a forgetful tensor functor
$\omega:\Rep_{\mathcal D}(G,\phi)\to\mathcal D$.
\end{definition}

\begin{theorem}[Relative Tannaka duality]
The assignments
\[
 (\mathcal C,F)\longmapsto
 (\underline{\Aut}^{\otimes}(F),\psi_F),
 \qquad
 (G,\phi)\longmapsto
 (\Rep_{\mathcal D}(G,\phi),\omega)
\]
give mutually inverse \emph{contravariant} equivalences between
$\operatorname{Symm}(\mathcal D)$ and the category of $\mathcal D$-groups.
\end{theorem}

This is the relative reconstruction theorem; see
\cite[\S4, Theorem 4.2.1]{CEOIncompressible}. Its full reconstruction
proof is not part of this lecture. The variance can be checked directly:
a pair $(T,\alpha)$ as above sends an automorphism $\xi$ of $F_2$ to
\[
 \alpha^{-1}(\xi T)\alpha
       \in\underline{\Aut}^{\otimes}(F_1).
\]
Thus a functor between representation categories produces a group
homomorphism in the opposite direction. For $\mathcal D=\Vect_k$,
the fundamental group is trivial and we recover ordinary neutral
Tannaka duality.

\begin{example}
For $\mathcal D=\SVec_k$, $\operatorname{char}k\ne2$, an internal
group is an affine supergroup. A compatible map $\mu_2\to G$ is
specified by an even element $z\in G(k)$ with $z^2=1$ whose conjugation
action is the parity automorphism of $\mathcal O(G)$. The category
$\Rep(G,z)$ consists of representations on which $z$ acts by parity.
If $G$ is purely even, the condition makes $z$ central. For a general
supergroup, centrality must not replace the parity-conjugation condition.
\end{example}

\subsection{The characteristic-zero classification}

Now suppose that $k$ is algebraically closed of characteristic zero.
Deligne's theorem states that a symmetric tensor category admits a
fiber functor to $\SVec_k$ if and only if it has moderate growth.
Consequently it is of the form $\Rep(G,z)$ above. The fiber functor,
when it exists, is unique up to monoidal natural isomorphism; the
isomorphism need not be unique. These statements apply to arbitrary
essentially small pretannakian categories; no tensor-generation
hypothesis is imposed. For this formulation and uniqueness, see
\cite[\S\S1.1, 1.5]{CEOIncompressible}. In the ordinary Tannakian case,
uniqueness without a finite-generation hypothesis also follows from
\cite[Corollary 1.2]{WibmerTorsors}.

The ordinary Tannakian criterion is stronger: there is a fiber functor
to $\Vect_k$ if and only if, for every object $X$, some exterior power
$\bigwedge^{n_X}X$ vanishes. The bound may depend on $X$. The next
lecture explains why arbitrary Schur functors replace exterior powers
in the super version of this criterion.
